\section{Basics \& Essential Snippets}

\subsection{Template Setup \& I/O Optimizations}
\textbf{Note:} the bits library is a GCC-specific header that includes almost every standard library in C++. It slows down compilation slightly but saves time typing during an exam.

\begin{lstlisting}
#include <bits/stdc++.h> 
using namespace std;

// Replaces all 'int' with 'long long' to prevent integer overflow globally.
// WARNING: If you use this, your main function MUST be declared as int32_t.

#define int long long

signed main() {
    // I/O Optimizations: Disconnects C and C++ standard streams.
    // Makes cin/cout almost as fast as scanf/printf.
    ios_base::sync_with_stdio(false);
    cin.tie(NULL);

    int num_elements, target_value;
    
    // Reading until End of File (EOF). 
    // Useful when the problem doesn't specify the number of test cases.
    while (cin >> num_elements >> target_value) {
        // Process each test case here
    }
    
    return 0;
}
\end{lstlisting}

\subsection{Basic Structures, Vectors, and Matrices}
\begin{lstlisting}
#include <iostream>
#include <vector>
#include <algorithm>

using namespace std;

// 1. Defining a Struct
struct Player {
    int score;
    int penalties;
};

// 2. Comparing Two Elements (Custom Sorting)
// Returns true if 'a' should be placed BEFORE 'b' in the sorted array
bool comparePlayers(const Player& a, const Player& b) {
    if (a.score != b.score) {
        return a.score > b.score; // Primary rule: Sort descending by score
    }
    return a.penalties < b.penalties; // Tie-breaker: Sort ascending by penalties
}

void standardDataStructures() {
    int num_items = 10;
    int rows = 5, cols = 5;

    // 3. Defining a Vector (1D Array)
    // Space Complexity: O(N)
    // Creates a vector of size 10, with all elements initialized to 0
    vector<int> one_dimensional_array(num_items, 0); 

    // 4. Defining a Matrix (2D Vector)
    // Space Complexity: O(Rows * Cols)
    // Creates a 5x5 matrix, with all elements initialized to -1
    vector<vector<int>> matrix(rows, vector<int>(cols, -1));

    // 5. Selecting Maximum Element
    vector<int> numbers = {1, 5, 3, 9, 2};
    
    // Time Complexity: O(N)
    // Note: max_element returns an ITERATOR. We use '*' to get the actual value.
    int max_value = *max_element(numbers.begin(), numbers.end());
}
\end{lstlisting}

\subsection{Backtracking to Find DP Solutions}
\textbf{Concept:} Once a DP table is filled, you start from the final optimal state (usually the bottom-right corner of the matrix) and trace backward. You compare the current state to previous states to determine which ``choice'' led you to the current optimal value.

\begin{lstlisting}
#include <iostream>
#include <vector>

using namespace std;

// Example: Recovering the items chosen in a 0/1 Knapsack Problem
// dp_table[i][w] stores the max value using the first 'i' items with capacity 'w'.
vector<int> getChosenItems(const vector<vector<int>>& dp_table, const vector<int>& weights, int num_items, int max_capacity) {
    // Space Complexity: O(N) to store the result
    vector<int> chosen_items;
    
    int current_weight = max_capacity;

    // Time Complexity: O(N) - We backtrack step-by-step through the items
    for (int i = num_items; i > 0; --i) {
        
        // If the value in the DP table changed compared to the row above it,
        // it means we MUST have included item 'i' to achieve this value.
        if (dp_table[i][current_weight] != dp_table[i-1][current_weight]) {
            
            // Note: DP table items are 1-indexed, but array weights are 0-indexed.
            int original_item_index = i - 1;
            chosen_items.push_back(original_item_index);
            
            // Deduct the weight of the chosen item to move left in the DP table
            current_weight -= weights[original_item_index];
        }
    }
    
    // Optional: reverse(chosen_items.begin(), chosen_items.end()); 
    // Do this if you need the items in the order they originally appeared.
    return chosen_items;
}
\end{lstlisting}